\documentclass[../main.tex]{subfiles}
\begin{document}

\chapter{オペレーティングシステム入門}
\lessoninfo{
  video={P1L2},
  textbook={OSTEP \cite{ostep} ch.~2，6；Silberschatz ら \cite{silberschatz2018} ch.~1--2},
  keywords={オペレーティングシステム，抽象化，調停，機構，方針，カーネルモード，ユーザモード，トラップ，システムコール，シグナル，モノリシックOS，モジュール型OS，マイクロカーネル，Linux},
}

あらゆるアプリケーションは，実行されるマシンの CPU，メモリ，デバイスを
使うためにオペレーティングシステムに依存している．この章では，
オペレーティングシステムとは何か，どのような役割を果たすかを定義し，
それを構成する要素---抽象化，機構，方針---と，その設計を導く原則を導入
する．次に，アプリケーションとオペレーティングシステムの間の保護境界と，
アプリケーションがシステムコールによってそれを越える仕組みを調べ，
オペレーティングシステムの構成の仕方を Linux と Mac OS X を例に比較
する．ここで導入するプロセス，スレッド，メモリページ，ファイルといった
抽象化は，以降の章の主題である．

\section{オペレーティングシステムとは何か}
\label{sec:l01-what}

計算機システムは，ハードウェア---1つ以上の CPU（それぞれ複数のコアを
持ちうる），主記憶，ハードディスクや SSD，USB フラッシュドライブなどの
記憶装置，Ethernet や WiFi のためのネットワークインタフェース，
グラフィックスプロセッサ（GPU）---と，それを使うアプリケーションから
なる．\source{P1L2，オペレーティングシステムとは何か，オペレーティング
システムの定義；OSTEP ch.~2；Silberschatz ら ch.~1．}%
通常は複数のアプリケーションが同時に実行され，そのすべてがこれらの
ハードウェア構成要素を必要とする．オペレーティングシステムは，両者の間に
位置するソフトウェアの層である（\cref{fig:l01-os-layers}）．

\begin{definition}[オペレーティングシステム]
  \term[おぺれーてぃんぐしすてむ]{オペレーティングシステム}[operating system]%
  （OS）とは，次の性質を持つシステムソフトウェアの層である．
  \begin{itemize}
    \item 下位のハードウェアに直接かつ特権的にアクセスできる．
    \item ハードウェアの複雑さをアプリケーションから隠す．
    \item あらかじめ定められた方針に従って，1つ以上のアプリケーションに
      代わってハードウェアを管理する．
    \item アプリケーションが互いに隔離され保護されることを保証する．
  \end{itemize}
\end{definition}
\index{OS|see{オペレーティングシステム}}

\begin{figure}[htbp]
  \centering
  % The operating system as the layer between applications and hardware.
% Usage: % The operating system as the layer between applications and hardware.
% Usage: % The operating system as the layer between applications and hardware.
% Usage: \input{figures/l01-os-layers}   (width ~110 mm)
\begin{tikzpicture}[x=1mm, y=1mm, font=\footnotesize\sffamily,
  app/.style={draw, rounded corners=2pt, fill=white, minimum height=7mm,
              minimum width=30mm, inner sep=1pt},
  comp/.style={draw, fill=white, minimum height=7mm, minimum width=23mm,
               inner sep=1pt, align=center, font=\scriptsize\sffamily},
  hw/.style={draw, fill=ln-box, minimum height=8mm, minimum width=16mm,
             inner sep=1pt, align=center, font=\scriptsize\sffamily},
  lab/.style={font=\scriptsize\sffamily, text=ln-muted},
  >={Stealth[length=1.6mm]}]
  % ---- applications
  \foreach \x in {17,55,93} {
    \node[app] at (\x,44) {\iflangja{アプリケーション}{application}};
    \draw[<->] (\x,40.5) -- (\x,32);
  }
  \node[lab, anchor=west] at (94.5,36.2) {\iflangja{システムコール}{system calls}};
  % ---- operating system
  \draw[thick, draw=ln-accent, fill=ln-box, rounded corners=2pt]
    (0,14) rectangle (110,32);
  \node[font=\footnotesize\sffamily\bfseries, text=ln-accent, anchor=north west]
    at (1,31.5) {\iflangja{オペレーティングシステム}{operating system}};
  \node[comp] at (14,20.5)   {\iflangja{スケジューラ}{scheduler}};
  \node[comp] at (41.3,20.5) {\iflangja{メモリマネージャ}{memory manager}};
  \node[comp] at (68.6,20.5) {\iflangja{ファイルシステム}{file systems}};
  \node[comp] at (96,20.5)   {\iflangja{デバイスドライバ}{device drivers}};
  % ---- hardware
  \foreach \x/\t in {8/{CPU},
                     26.8/{\iflangja{メモリ}{memory}},
                     45.6/{GPU},
                     64.4/{\iflangja{ネットワーク\\カード}{network\\card}},
                     83.2/{\iflangja{ディスク，\\SSD}{disk,\\SSD}},
                     102/{\iflangja{USB\\デバイス}{USB\\devices}}} {
    \node[hw] at (\x,4) {\t};
    \draw[<->] (\x,8) -- (\x,14);
  }
\end{tikzpicture}
   (width ~110 mm)
\begin{tikzpicture}[x=1mm, y=1mm, font=\footnotesize\sffamily,
  app/.style={draw, rounded corners=2pt, fill=white, minimum height=7mm,
              minimum width=30mm, inner sep=1pt},
  comp/.style={draw, fill=white, minimum height=7mm, minimum width=23mm,
               inner sep=1pt, align=center, font=\scriptsize\sffamily},
  hw/.style={draw, fill=ln-box, minimum height=8mm, minimum width=16mm,
             inner sep=1pt, align=center, font=\scriptsize\sffamily},
  lab/.style={font=\scriptsize\sffamily, text=ln-muted},
  >={Stealth[length=1.6mm]}]
  % ---- applications
  \foreach \x in {17,55,93} {
    \node[app] at (\x,44) {\iflangja{アプリケーション}{application}};
    \draw[<->] (\x,40.5) -- (\x,32);
  }
  \node[lab, anchor=west] at (94.5,36.2) {\iflangja{システムコール}{system calls}};
  % ---- operating system
  \draw[thick, draw=ln-accent, fill=ln-box, rounded corners=2pt]
    (0,14) rectangle (110,32);
  \node[font=\footnotesize\sffamily\bfseries, text=ln-accent, anchor=north west]
    at (1,31.5) {\iflangja{オペレーティングシステム}{operating system}};
  \node[comp] at (14,20.5)   {\iflangja{スケジューラ}{scheduler}};
  \node[comp] at (41.3,20.5) {\iflangja{メモリマネージャ}{memory manager}};
  \node[comp] at (68.6,20.5) {\iflangja{ファイルシステム}{file systems}};
  \node[comp] at (96,20.5)   {\iflangja{デバイスドライバ}{device drivers}};
  % ---- hardware
  \foreach \x/\t in {8/{CPU},
                     26.8/{\iflangja{メモリ}{memory}},
                     45.6/{GPU},
                     64.4/{\iflangja{ネットワーク\\カード}{network\\card}},
                     83.2/{\iflangja{ディスク，\\SSD}{disk,\\SSD}},
                     102/{\iflangja{USB\\デバイス}{USB\\devices}}} {
    \node[hw] at (\x,4) {\t};
    \draw[<->] (\x,8) -- (\x,14);
  }
\end{tikzpicture}
   (width ~110 mm)
\begin{tikzpicture}[x=1mm, y=1mm, font=\footnotesize\sffamily,
  app/.style={draw, rounded corners=2pt, fill=white, minimum height=7mm,
              minimum width=30mm, inner sep=1pt},
  comp/.style={draw, fill=white, minimum height=7mm, minimum width=23mm,
               inner sep=1pt, align=center, font=\scriptsize\sffamily},
  hw/.style={draw, fill=ln-box, minimum height=8mm, minimum width=16mm,
             inner sep=1pt, align=center, font=\scriptsize\sffamily},
  lab/.style={font=\scriptsize\sffamily, text=ln-muted},
  >={Stealth[length=1.6mm]}]
  % ---- applications
  \foreach \x in {17,55,93} {
    \node[app] at (\x,44) {\iflangja{アプリケーション}{application}};
    \draw[<->] (\x,40.5) -- (\x,32);
  }
  \node[lab, anchor=west] at (94.5,36.2) {\iflangja{システムコール}{system calls}};
  % ---- operating system
  \draw[thick, draw=ln-accent, fill=ln-box, rounded corners=2pt]
    (0,14) rectangle (110,32);
  \node[font=\footnotesize\sffamily\bfseries, text=ln-accent, anchor=north west]
    at (1,31.5) {\iflangja{オペレーティングシステム}{operating system}};
  \node[comp] at (14,20.5)   {\iflangja{スケジューラ}{scheduler}};
  \node[comp] at (41.3,20.5) {\iflangja{メモリマネージャ}{memory manager}};
  \node[comp] at (68.6,20.5) {\iflangja{ファイルシステム}{file systems}};
  \node[comp] at (96,20.5)   {\iflangja{デバイスドライバ}{device drivers}};
  % ---- hardware
  \foreach \x/\t in {8/{CPU},
                     26.8/{\iflangja{メモリ}{memory}},
                     45.6/{GPU},
                     64.4/{\iflangja{ネットワーク\\カード}{network\\card}},
                     83.2/{\iflangja{ディスク，\\SSD}{disk,\\SSD}},
                     102/{\iflangja{USB\\デバイス}{USB\\devices}}} {
    \node[hw] at (\x,4) {\t};
    \draw[<->] (\x,8) -- (\x,14);
  }
\end{tikzpicture}

  \caption{アプリケーションとハードウェアの間のオペレーティングシステム．
    アプリケーションはシステムコールを通じてサービスを要求し，
    スケジューラ，メモリマネージャ，ファイルシステム，デバイスドライバなどの
    構成要素がそれをハードウェア上で実行する．}
  \label{fig:l01-os-layers}
\end{figure}

要するに，オペレーティングシステムはハードウェアを\emph{抽象化}して
アプリケーションから見た姿を単純にし，その利用を\emph{調停}して，
アプリケーションがハードウェアをどのように共有するかを監督・制御する．
この定義は3つの役割を挙げている．
\begin{description}
  \item[ハードウェアの複雑さの隠蔽] オペレーティングシステムは，
    \term[ちゅうしょうか]{抽象化}[abstraction]，すなわち背後にある
    デバイスの詳細を隠す，より単純で高水準なインタフェースを通じて
    ハードウェアをアプリケーションに提示する．データを保存する
    アプリケーションは，特定のディスクのセクタやブロックを指定しない．
    \emph{ファイル}に書き込み，オペレーティングシステムがファイル操作を，
    実際に備わっている記憶装置に対する操作へ変換する．同じように
    \emph{ソケット}は，ネットワークインタフェースによらずネットワーク
    通信を表す．デバイスが変わっても抽象化は変わらないので，
    アプリケーションは異なるハードウェア上でも変更なしに動作する．
  \item[資源管理] オペレーティングシステムは
    \term[ちょうてい]{調停}[arbitration]によってハードウェア資源を管理
    する．すなわち，どのアプリケーションがどの資源をどれだけ使うかを
    決め，アプリケーションにメモリを割り当て，それらを CPU 上に
    スケジュールし，デバイスへのアクセスを制御する．その判断は方針に
    従う．例えば，資源へのアクセスが公平であること，あるいは1つの
    アプリケーションが使う資源に上限を設けることである．
  \item[隔離と保護] 同時に実行されるアプリケーションは互いに干渉して
    はならない．オペレーティングシステムは，各アプリケーションに資源の
    取り分---例えば物理メモリの自分の領域---を与えることで
    \term[かくり]{隔離}[isolation]を提供し，その実行が他の
    アプリケーションに影響されないようにする．また，割り当てられて
    いない資源，例えば他のアプリケーションやオペレーティングシステム
    自身のメモリへのアクセスを禁じることで\term[ほご]{保護}[protection]%
    を提供する．
\end{description}

したがってオペレーティングシステムは，CPU スケジューラ，メモリマネージャ，
デバイスドライバ，ファイルシステムといった構成要素からなる．表計算ソフトや
メディアプレーヤのように，これらのサービスを利用するだけの
アプリケーションプログラムは，たとえ一緒に配布されていても，その一部
ではない．

\begin{example}
  オペレーティングシステムのある機能は，デバイスへのインタフェースを
  単純にするときには抽象化であり，資源をアプリケーション間で分けるとき
  には調停である．プログラムが，プリンタの機種の命令言語を知らずに
  文書をプリンタへ送れるようにすることは抽象化であり，同時に印刷する
  3つのプログラムの文書がどの順でプリンタに届くかを決めることは調停で
  ある．同様に，Ethernet でも WiFi でも同じソケットインタフェースを
  提供することは抽象化であり，1つのアプリケーションが消費してよい
  ネットワーク帯域を制限することは調停である．
\end{example}

\subsection{オペレーティングシステムの例}

オペレーティングシステムは，対象とする環境によって異なる．\source{P1L2，
オペレーティングシステムの例．}デスクトップやサーバのマシンでは，主な例は
Microsoft Windows と UNIX 系のシステムであり，後者には Mac OS X
（BSD UNIX に由来する）や Linux が含まれる．スマートフォンなどの組込みデバイス
では，Linux に基づく Android，iOS，Symbian などがある．いずれもその環境
に適した設計上の選択を行っている．以降の章では主に Linux を例として
用いる．

\section{抽象化，機構，方針}
\label{sec:l01-elements}

オペレーティングシステムは3種類の要素から構築される．\source{P1L2，OS の
構成要素，OS の構成要素の例；Silberschatz ら ch.~2．}その抽象化は，実行中
のアプリケーション---プロセスとスレッド---と，ハードウェア資源---ファイル，
ソケット，メモリページ---を表す．オペレーティングシステムはこれらの抽象化
に対する操作を提供し，その操作をどのように適用するかを決める．

\begin{definition}[機構と方針]
  \term[きこう]{機構}[mechanism]とは，プロセスの生成やスケジューリング，
  ファイルのオープンや書き込み，メモリの割り当てなど，オペレーティング
  システムがその抽象化に対して提供する操作である．
  \term[ほうしん]{方針}[policy]とは，メモリに対する LRU（least recently
  used）やスケジューリングに対する EDF（earliest deadline first）の
  ように，ハードウェアを管理するために機構をどう使うかを定める規則で
  ある．
\end{definition}

メモリ管理は，3つの要素がどのように組み合わさるかを示している
（\cref{tab:l01-elements}）．抽象化はメモリページ，すなわち固定長の
メモリブロックである．機構はページを割り当て，それをプロセスのアドレス
空間に対応付けて，プロセスがそれを使えるようにする．物理メモリ（DRAM）は
通常すべてのプロセスの要求の合計より小さいので，オペレーティングシステム
は一部のページをディスクへ移し，再びアクセスされたときに戻す．
どのページを移すかを決めるのが方針である．LRU では，最も長くアクセス
されていないページを，近いうちに必要になる可能性が最も低いという仮定の
もとで移す．

\begin{table}[htbp]
  \centering
  \caption{オペレーティングシステムの要素．一般の場合とメモリ管理の場合．}
  \label{tab:l01-elements}
  \small
  \begin{tabular}{@{}l>{\raggedright}p{47mm}>{\raggedright\arraybackslash}p{45mm}@{}}
    \toprule
    要素 & 例 & メモリ管理 \\
    \midrule
    抽象化 & プロセス，スレッド，ファイル，ソケット，メモリページ & メモリページ \\
    機構 & 生成，スケジュール，オープン，書き込み，割り当て
      & ページを割り当て，プロセスに対応付ける \\
    方針 & LRU，EDF & LRU: 最も長くアクセスされていないページをディスクへ移す \\
    \bottomrule
  \end{tabular}
\end{table}

\section{設計原則}
\label{sec:l01-principles}

これらの要素の設計を導く原則が2つある．\source{P1L2，OS の設計原則；
Silberschatz ら ch.~2．}第一は
\term[きこうとほうしんのぶんり]{機構と方針の分離}[separation of mechanism and policy]%
である．オペレーティングシステムは多くの方針を支えられる柔軟な機構を
実装し，方針は機構を変更することなく選択したり置き換えたりできる．
ディスクへ移すページの選択では，ページアクセスの記録，ページのディスクへの
移動，その対応付けの更新という同じ機構が，LRU，LFU（least frequently
used），さらにはランダム選択も支える．どれが最も適しているかは，システムが
使われる状況に依存する．

\begin{example}
  物理メモリが4ページを保持しており，5つ目のページが必要になったので，
  4つのうち1つをディスクへ移さなければならない．時刻
  \SI{160}{\milli\second} における記録されたアクセスは次のとおりである．
  \begin{center}
    \small
    \begin{tabular}{@{}lcccc@{}}
      \toprule
      ページ & A & B & C & D \\
      \midrule
      最終アクセス（\si{\milli\second}） & 120 & 40 & 95 & 150 \\
      アクセス回数 & 9 & 25 & 3 & 12 \\
      \bottomrule
    \end{tabular}
  \end{center}
  LRU は最終アクセスが最も過去にあるページ B を，LFU はアクセス回数が
  最も少ないページ C を移し，ランダム選択は各ページを確率 $1/4$ で移す．
  いずれの場合も実行される機構は同じであり，選ばれたページがディスクへ
  書き出され，その対応付けが無効化される．
\end{example}

第二の原則は\term[ひんぱんにおこるばあいのさいてきか]{頻繁に起こる場合の最適化}[optimize for the common case]%
である．設計者は，オペレーティングシステムがどこで使われるか，利用者が
そのマシンでどのアプリケーションを実行するか，そのワークロードの要件が
何かを理解しなければならない．その上で，この頻繁に起こる場合に
適し，かつ下位の抽象化と機構で支えられる方針を選ぶ．例えば主に
データベースを実行するサーバは多数のランダムな読み書きを行うのに対し，
メディアサーバは大きなファイルを逐次に読む．一方のワークロード
に適した方針は，他方にはうまく合わない．

\begin{keypoint}
  抽象化は複雑さを単純にし，機構は抽象化に対して操作を行い，方針は機構を
  どう使うかを決める．機構を一般的に保つことで，同じオペレーティングシステムが，
  その頻繁に起こる場合に合った方針を採用できる．
\end{keypoint}

\section{ユーザ/カーネル保護境界}
\label{sec:l01-boundary}

オペレーティングシステムはハードウェアに直接アクセスする必要があるので，
アプリケーションが持たない特権をもって実行されなければならない．
\source{P1L2，ユーザ/カーネル保護境界；OSTEP ch.~6；Silberschatz ら
ch.~1．}%
そこで計算機プラットフォームは，少なくとも2つの実行モードを区別する．

\begin{definition}[カーネルモードとユーザモード]
  特権モードである\term[かーねるもーど]{カーネルモード}[kernel mode]では，
  ソフトウェアはプロセッサのすべての命令を実行でき，ハードウェアに直接
  アクセスできる．非特権モードである
  \term[ゆーざもーど]{ユーザモード}[user mode]では，ハードウェアに直接
  アクセスする命令は禁じられる．
\end{definition}

カーネルモードで実行されるオペレーティングシステムの部分が，その
\term[かーねる]{カーネル}[kernel]である．アプリケーションはユーザモードで
実行される．2つのモードの区別はハードウェアによって支えられる．CPU は
現在のモードを記録する\term[もーどびっと]{モードビット}[mode bit]を持ち，
ハードウェアを直接操作する命令，すなわち
\term[とっけんめいれい]{特権命令}[privileged instruction]は，モードビット
がカーネルモードを示す場合にのみ実行される．ユーザモードで特権命令を実行
しようとしても効果はなく，トラップが発生する．

\begin{definition}[トラップ]
  \term[とらっぷ]{トラップ}[trap]とは，実行中のアプリケーションから
  オペレーティングシステムへの制御の移行であり，ハードウェアによって
  行われる．アプリケーションは中断され，CPU はカーネルモードに切り替わり，
  オペレーティングシステムがあらかじめ登録しておいたカーネル内の位置から
  実行が続く．
\end{definition}

特権命令や，特権を必要とするメモリアクセスによるトラップの後，
オペレーティングシステムはトラップの原因が何か，そのアクセスを許すべきか
を判定する．\margin{x86 プロセッサは4つの特権レベル（第12章の保護リング）
を備え，OS はうち2つ以上を使う．}アプリケーションが不正なことを試みた場合
には，それを終了させることもある．アプリケーションとオペレーティングシステム
は3つの方法でやり取りする（\cref{fig:l01-user-kernel}）．
\begin{description}
  \item[トラップ] アプリケーションが特権を必要とする操作を行ったときに，
    制御をオペレーティングシステムへ移す．
  \item[システムコール] 明示的に要求される．
    \term[しすてむこーる]{システムコール}[system call]とは，
    アプリケーションが自分に代わってサービスや特権的なアクセスを行う
    ようオペレーティングシステムに依頼できるように，オペレーティング
    システムが公開している操作である．例えばファイルに対する
    \texttt{open}，ソケットに対する \texttt{send}，メモリに対する
    \texttt{mmap} がある．これらの操作の集合が\emph{システムコール
    インタフェース}である．
  \item[シグナル] 逆向きに伝わる．\term[しぐなる]{シグナル}[signal]とは，
    オペレーティングシステムがアプリケーションに渡す通知であり，例えば
    アプリケーションが設定したタイマが満了したことを伝える．
\end{description}

\begin{figure}[htbp]
  \centering
  % The user/kernel protection boundary and the ways to cross it.
% Usage: % The user/kernel protection boundary and the ways to cross it.
% Usage: % The user/kernel protection boundary and the ways to cross it.
% Usage: \input{figures/l01-user-kernel}   (width ~112 mm)
\begin{tikzpicture}[x=1mm, y=1mm, font=\footnotesize\sffamily,
  app/.style={draw, rounded corners=2pt, fill=white, minimum height=8mm,
              minimum width=30mm, inner sep=1pt},
  big/.style={draw, rounded corners=2pt, minimum height=8mm,
              minimum width=112mm, inner sep=1pt},
  lab/.style={font=\scriptsize\sffamily, align=center},
  lvl/.style={font=\scriptsize\sffamily, text=ln-muted, align=right},
  >={Stealth[length=1.8mm]}]
  \node[app] (a) at (28,42) {\iflangja{アプリケーション}{application}};
  \node[app] (b) at (78,42) {\iflangja{アプリケーション}{application}};
  \draw[dashed, ln-rule, thick] (0,26) -- (112,26);
  \node[lvl, anchor=south east] at (112,26.8) {\iflangja{ユーザレベル（非特権）}{user level (unprivileged)}};
  \node[lvl, anchor=north east] at (112,25.2) {\iflangja{カーネルレベル（特権）}{kernel level (privileged)}};
  \node[big, fill=ln-box, thick, draw=ln-accent] (k) at (56,15)
    {\iflangja{OS カーネル}{OS kernel}};
  \node[big, fill=white] (h) at (56,0) {\iflangja{ハードウェア}{hardware}};
  % system call and signal
  \draw[->, thick, ln-accent] (22,38) -- (22,19);
  \node[lab, anchor=east] at (21,32) {\iflangja{システムコール}{system call}};
  \draw[->, thick, ln-accent] (34,19) -- (34,38);
  \node[lab, anchor=west] at (35,32) {\iflangja{シグナル}{signal}};
  % trap
  \draw[->, thick, ln-caution] (78,38) -- (78,19);
  \node[lab, anchor=east] at (77,32)
    {\iflangja{トラップ\\（特権操作）}{trap (privileged\\operation)}};
  % privileged access
  \draw[<->, thick] (56,11) -- (56,4);
  \node[lab, anchor=west] at (57,7.5)
    {\iflangja{ハードウェアへの直接アクセス}{direct hardware access}};
\end{tikzpicture}
   (width ~112 mm)
\begin{tikzpicture}[x=1mm, y=1mm, font=\footnotesize\sffamily,
  app/.style={draw, rounded corners=2pt, fill=white, minimum height=8mm,
              minimum width=30mm, inner sep=1pt},
  big/.style={draw, rounded corners=2pt, minimum height=8mm,
              minimum width=112mm, inner sep=1pt},
  lab/.style={font=\scriptsize\sffamily, align=center},
  lvl/.style={font=\scriptsize\sffamily, text=ln-muted, align=right},
  >={Stealth[length=1.8mm]}]
  \node[app] (a) at (28,42) {\iflangja{アプリケーション}{application}};
  \node[app] (b) at (78,42) {\iflangja{アプリケーション}{application}};
  \draw[dashed, ln-rule, thick] (0,26) -- (112,26);
  \node[lvl, anchor=south east] at (112,26.8) {\iflangja{ユーザレベル（非特権）}{user level (unprivileged)}};
  \node[lvl, anchor=north east] at (112,25.2) {\iflangja{カーネルレベル（特権）}{kernel level (privileged)}};
  \node[big, fill=ln-box, thick, draw=ln-accent] (k) at (56,15)
    {\iflangja{OS カーネル}{OS kernel}};
  \node[big, fill=white] (h) at (56,0) {\iflangja{ハードウェア}{hardware}};
  % system call and signal
  \draw[->, thick, ln-accent] (22,38) -- (22,19);
  \node[lab, anchor=east] at (21,32) {\iflangja{システムコール}{system call}};
  \draw[->, thick, ln-accent] (34,19) -- (34,38);
  \node[lab, anchor=west] at (35,32) {\iflangja{シグナル}{signal}};
  % trap
  \draw[->, thick, ln-caution] (78,38) -- (78,19);
  \node[lab, anchor=east] at (77,32)
    {\iflangja{トラップ\\（特権操作）}{trap (privileged\\operation)}};
  % privileged access
  \draw[<->, thick] (56,11) -- (56,4);
  \node[lab, anchor=west] at (57,7.5)
    {\iflangja{ハードウェアへの直接アクセス}{direct hardware access}};
\end{tikzpicture}
   (width ~112 mm)
\begin{tikzpicture}[x=1mm, y=1mm, font=\footnotesize\sffamily,
  app/.style={draw, rounded corners=2pt, fill=white, minimum height=8mm,
              minimum width=30mm, inner sep=1pt},
  big/.style={draw, rounded corners=2pt, minimum height=8mm,
              minimum width=112mm, inner sep=1pt},
  lab/.style={font=\scriptsize\sffamily, align=center},
  lvl/.style={font=\scriptsize\sffamily, text=ln-muted, align=right},
  >={Stealth[length=1.8mm]}]
  \node[app] (a) at (28,42) {\iflangja{アプリケーション}{application}};
  \node[app] (b) at (78,42) {\iflangja{アプリケーション}{application}};
  \draw[dashed, ln-rule, thick] (0,26) -- (112,26);
  \node[lvl, anchor=south east] at (112,26.8) {\iflangja{ユーザレベル（非特権）}{user level (unprivileged)}};
  \node[lvl, anchor=north east] at (112,25.2) {\iflangja{カーネルレベル（特権）}{kernel level (privileged)}};
  \node[big, fill=ln-box, thick, draw=ln-accent] (k) at (56,15)
    {\iflangja{OS カーネル}{OS kernel}};
  \node[big, fill=white] (h) at (56,0) {\iflangja{ハードウェア}{hardware}};
  % system call and signal
  \draw[->, thick, ln-accent] (22,38) -- (22,19);
  \node[lab, anchor=east] at (21,32) {\iflangja{システムコール}{system call}};
  \draw[->, thick, ln-accent] (34,19) -- (34,38);
  \node[lab, anchor=west] at (35,32) {\iflangja{シグナル}{signal}};
  % trap
  \draw[->, thick, ln-caution] (78,38) -- (78,19);
  \node[lab, anchor=east] at (77,32)
    {\iflangja{トラップ\\（特権操作）}{trap (privileged\\operation)}};
  % privileged access
  \draw[<->, thick] (56,11) -- (56,4);
  \node[lab, anchor=west] at (57,7.5)
    {\iflangja{ハードウェアへの直接アクセス}{direct hardware access}};
\end{tikzpicture}

  \caption{ユーザ/カーネル保護境界．ハードウェアにアクセスするのは
    カーネルだけであり，アプリケーションはシステムコールまたはトラップに
    よってカーネルに入り，カーネルはシグナルでアプリケーションに通知
    する．}
  \label{fig:l01-user-kernel}
\end{figure}

\section{システムコール}
\label{sec:l01-syscalls}

\cref{fig:l01-syscall-flow} は，ユーザプロセスがシステムコールを実行する
ときの一連の出来事を示している．\source{P1L2，システムコールのフロー
チャート；OSTEP ch.~6；Silberschatz ら ch.~1--2．}プロセスは，
オペレーティングシステムのサービスが必要になってシステムコールを呼び出す
までユーザモードで実行される．この呼び出しはカーネルへトラップする．
モードビットがカーネルモードに切り替わり，実行コンテキストがユーザ
プロセスのものからカーネルのものへ変わり，引数が渡され，そのコールを
実装するカーネルのコードへ実行が移る．カーネルはカーネルモードでコールを
実行する．完了すると，モードビットがユーザモードに戻され，プロセスの実行
コンテキストが復元され，結果が返され，プロセスはコールの次の命令から実行
を続ける．

\begin{figure}[htbp]
  \centering
  % Sequence of a system call: transition from user mode to kernel mode and back.
% Usage: % Sequence of a system call: transition from user mode to kernel mode and back.
% Usage: % Sequence of a system call: transition from user mode to kernel mode and back.
% Usage: \input{figures/l01-syscall-flow}   (width ~118 mm)
\begin{tikzpicture}[x=1mm, y=1mm, font=\footnotesize\sffamily,
  st/.style={draw, rounded corners=2pt, align=center, minimum height=9mm,
             inner sep=1pt, font=\scriptsize\sffamily},
  lab/.style={font=\scriptsize\sffamily, align=center},
  lvl/.style={font=\scriptsize\sffamily\bfseries, text=ln-muted},
  >={Stealth[length=1.8mm]}]
  % kernel-mode region below the dashed boundary
  \fill[ln-box] (0,-6) rectangle (118,15);
  \draw[dashed, ln-rule, thick] (0,15) -- (118,15);
  \node[lvl, anchor=north west] at (0,32)
    {\iflangja{ユーザモード（モードビット = 1）}{user mode (mode bit = 1)}};
  \node[lvl, anchor=south west] at (0,-5.5)
    {\iflangja{カーネルモード（モードビット = 0）}{kernel mode (mode bit = 0)}};
  \node[st, fill=white, minimum width=24mm, text width=22mm] (a) at (13,23)
    {\iflangja{ユーザプロセス\\実行中}{user process\\executing}};
  \node[st, fill=white, minimum width=22mm, text width=20mm] (b) at (43,23)
    {\iflangja{システムコール\\を呼び出す}{calls\\system call}};
  \node[st, fill=white, thick, draw=ln-accent, minimum width=24mm, text width=22mm]
    (c) at (68,4) {\iflangja{システムコール\\を実行}{execute\\system call}};
  \node[st, fill=white, minimum width=24mm, text width=22mm] (d) at (93,23)
    {\iflangja{システムコール\\から復帰}{return from\\system call}};
  \draw[->] (a) -- (b);
  \draw[->, thick, ln-accent] (b.south) |- (c.west);
  \node[lab, anchor=east] at (42,8.5)
    {\iflangja{トラップ\\モードビット = 0}{trap\\mode bit = 0}};
  \draw[->, thick, ln-accent] (c.east) -| (d.south);
  \node[lab, anchor=west] at (94,8.5)
    {\iflangja{復帰\\モードビット = 1}{return\\mode bit = 1}};
  \draw[->] (d.east) -- (118,23);
\end{tikzpicture}
   (width ~118 mm)
\begin{tikzpicture}[x=1mm, y=1mm, font=\footnotesize\sffamily,
  st/.style={draw, rounded corners=2pt, align=center, minimum height=9mm,
             inner sep=1pt, font=\scriptsize\sffamily},
  lab/.style={font=\scriptsize\sffamily, align=center},
  lvl/.style={font=\scriptsize\sffamily\bfseries, text=ln-muted},
  >={Stealth[length=1.8mm]}]
  % kernel-mode region below the dashed boundary
  \fill[ln-box] (0,-6) rectangle (118,15);
  \draw[dashed, ln-rule, thick] (0,15) -- (118,15);
  \node[lvl, anchor=north west] at (0,32)
    {\iflangja{ユーザモード（モードビット = 1）}{user mode (mode bit = 1)}};
  \node[lvl, anchor=south west] at (0,-5.5)
    {\iflangja{カーネルモード（モードビット = 0）}{kernel mode (mode bit = 0)}};
  \node[st, fill=white, minimum width=24mm, text width=22mm] (a) at (13,23)
    {\iflangja{ユーザプロセス\\実行中}{user process\\executing}};
  \node[st, fill=white, minimum width=22mm, text width=20mm] (b) at (43,23)
    {\iflangja{システムコール\\を呼び出す}{calls\\system call}};
  \node[st, fill=white, thick, draw=ln-accent, minimum width=24mm, text width=22mm]
    (c) at (68,4) {\iflangja{システムコール\\を実行}{execute\\system call}};
  \node[st, fill=white, minimum width=24mm, text width=22mm] (d) at (93,23)
    {\iflangja{システムコール\\から復帰}{return from\\system call}};
  \draw[->] (a) -- (b);
  \draw[->, thick, ln-accent] (b.south) |- (c.west);
  \node[lab, anchor=east] at (42,8.5)
    {\iflangja{トラップ\\モードビット = 0}{trap\\mode bit = 0}};
  \draw[->, thick, ln-accent] (c.east) -| (d.south);
  \node[lab, anchor=west] at (94,8.5)
    {\iflangja{復帰\\モードビット = 1}{return\\mode bit = 1}};
  \draw[->] (d.east) -- (118,23);
\end{tikzpicture}
   (width ~118 mm)
\begin{tikzpicture}[x=1mm, y=1mm, font=\footnotesize\sffamily,
  st/.style={draw, rounded corners=2pt, align=center, minimum height=9mm,
             inner sep=1pt, font=\scriptsize\sffamily},
  lab/.style={font=\scriptsize\sffamily, align=center},
  lvl/.style={font=\scriptsize\sffamily\bfseries, text=ln-muted},
  >={Stealth[length=1.8mm]}]
  % kernel-mode region below the dashed boundary
  \fill[ln-box] (0,-6) rectangle (118,15);
  \draw[dashed, ln-rule, thick] (0,15) -- (118,15);
  \node[lvl, anchor=north west] at (0,32)
    {\iflangja{ユーザモード（モードビット = 1）}{user mode (mode bit = 1)}};
  \node[lvl, anchor=south west] at (0,-5.5)
    {\iflangja{カーネルモード（モードビット = 0）}{kernel mode (mode bit = 0)}};
  \node[st, fill=white, minimum width=24mm, text width=22mm] (a) at (13,23)
    {\iflangja{ユーザプロセス\\実行中}{user process\\executing}};
  \node[st, fill=white, minimum width=22mm, text width=20mm] (b) at (43,23)
    {\iflangja{システムコール\\を呼び出す}{calls\\system call}};
  \node[st, fill=white, thick, draw=ln-accent, minimum width=24mm, text width=22mm]
    (c) at (68,4) {\iflangja{システムコール\\を実行}{execute\\system call}};
  \node[st, fill=white, minimum width=24mm, text width=22mm] (d) at (93,23)
    {\iflangja{システムコール\\から復帰}{return from\\system call}};
  \draw[->] (a) -- (b);
  \draw[->, thick, ln-accent] (b.south) |- (c.west);
  \node[lab, anchor=east] at (42,8.5)
    {\iflangja{トラップ\\モードビット = 0}{trap\\mode bit = 0}};
  \draw[->, thick, ln-accent] (c.east) -| (d.south);
  \node[lab, anchor=west] at (94,8.5)
    {\iflangja{復帰\\モードビット = 1}{return\\mode bit = 1}};
  \draw[->] (d.east) -- (118,23);
\end{tikzpicture}

  \caption{システムコールの実行．トラップによってカーネルモードに
    切り替わり，カーネルがコールを実行し，復帰によってユーザモードに
    戻る．モードビットは Silberschatz ら \cite{silberschatz2018} の
    慣習に従い，カーネルモードで 0，ユーザモードで 1 である．}
  \label{fig:l01-syscall-flow}
\end{figure}

システムコールを行うには，アプリケーションは引数とその他の必要なデータを
定められた場所に書き，\emph{システムコール番号}によってコールを呼び出す．
カーネルはこの番号から，どの操作が要求されているか，引数がいくつあるか，
どこにあるかを判断する．引数は，値として直接渡されるか，それらを保持する
メモリ領域のアドレスとして間接的に渡される．ここで述べる同期的な形式では，
プロセスはシステムコールが完了するまで待つ．非同期のシステムコールは後の
章で扱う．

\begin{example}
  次の C プログラムのどちらの呼び出しも，標準出力へ6バイトを書くよう
  Linux に依頼する．
  \begin{lstlisting}[language=C, numbers=none]
#define _GNU_SOURCE
#include <unistd.h>
#include <sys/syscall.h>

int main(void) {
  const char msg[] = "hello\n";
  /* システムコールの C ライブラリラッパ */
  write(1, msg, sizeof msg - 1);
  /* 同じシステムコールを番号で呼び出す */
  syscall(SYS_write, 1, msg, sizeof msg - 1);
  return 0;
}
\end{lstlisting}
  ライブラリ関数 \texttt{write} は，\texttt{write} システムコールの番号と
  その3つの引数をカーネルが期待する場所に置き，トラップを実行する．
  ファイルディスクリプタ~1（標準出力）と長さは直接渡され，文字列はその
  アドレス \texttt{msg} を通じて間接的に渡される．\margin{x86-64 Linux
  では，番号は \texttt{rax}，引数は \texttt{rdi}，\texttt{rsi}，
  \texttt{rdx} に置き，\texttt{syscall} 命令がトラップする．}
\end{example}

\subsection{ユーザ/カーネル境界の横断}

ユーザモードとカーネルモードの間の移行は，アプリケーションの実行に
欠かせない部分である．ファイルを読んだり，ネットワークでデータを送ったり，
メモリを割り当てたりするアプリケーションは，境界を越えなければならない．
\source{P1L2，OS 境界の横断．}この移行には3つの性質がある．
\begin{itemize}
  \item 不正な命令や特別な特権を必要とするメモリアクセスに対する
    トラップを通じて，ハードウェアに支えられている．
  \item 状態を保存・切り替えるために相当数の命令を実行する必要がある．
    Linux が動作する \SI{2}{\giga\hertz} のマシンでは，1回の移行に
    およそ \SIrange[range-phrase={--},range-units=single]{50}{100}{\nano\second}
    かかる．
  \item 局所性が切り替わる．カーネルはアプリケーションとは異なるコードを
    実行し異なるデータを使うので，その命令がプロセッサキャッシュ内の
    アプリケーションのデータを追い出す．アプリケーションが再開すると，
    データの一部を再び主記憶から取ってこなければならず，その分だけ遅く
    なる．
\end{itemize}
したがってユーザ/カーネル間の移行は安価ではない．\caution{システムコール
は C の手続き呼び出しのように見えるが，コストははるかに高い．}

\section{オペレーティングシステムのサービス}
\label{sec:l01-services}

オペレーティングシステムは，サービスを公開することでアプリケーションに
ハードウェアへのアクセスを与える．\source{P1L2，OS のサービス；
Silberschatz ら ch.~2．}最も基本的なレベルでは，サービスはハードウェア
構成要素に対応する．スケジューラは CPU へのアクセスを制御し，メモリマネージャは
同時に実行されるアプリケーションに物理メモリを割り当て，それらが互いの
メモリを上書きしないようにし，ブロックデバイスドライバはディスクなどの
ブロックデバイスへのアクセスを与える．その他のサービスは，1つの
ハードウェア構成要素に直接対応しない，より高水準の抽象化を提供する．
ファイルシステムのファイルやディレクトリがその例である．目的ごとに
まとめると，オペレーティングシステムはプロセス管理，ファイル管理，
デバイス管理，メモリ管理，ストレージ管理，セキュリティを提供する．

アプリケーションはこれらすべてのサービスにシステムコールを通じて到達
する．コールの名前や構成はオペレーティングシステムによって異なるが，
\cref{tab:l01-syscalls} の Windows と Unix のコールが示すように，
提供される機能は同等である．

\begin{table}[htbp]
  \centering
  \caption{Windows と Unix のシステムコールの例
    （Silberschatz ら \cite{silberschatz2018} による）．}
  \label{tab:l01-syscalls}
  \small
  \begin{tabular}{@{}>{\raggedright}p{24mm}>{\raggedright}p{55mm}%
      >{\raggedright\arraybackslash}p{32mm}@{}}
    \toprule
    分類 & Windows & Unix \\
    \midrule
    プロセス制御
      & \texttt{CreateProcess()}, \texttt{ExitProcess()},
        \texttt{WaitForSingleObject()}
      & \texttt{fork()}, \texttt{exit()}, \texttt{wait()} \\
    \addlinespace
    ファイル操作
      & \texttt{CreateFile()}, \texttt{ReadFile()}, \texttt{WriteFile()},
        \texttt{CloseHandle()}
      & \texttt{open()}, \texttt{read()}, \texttt{write()},
        \texttt{close()} \\
    \addlinespace
    デバイス操作
      & \texttt{SetConsoleMode()}, \texttt{ReadConsole()},
        \texttt{WriteConsole()}
      & \texttt{ioctl()}, \texttt{read()}, \texttt{write()} \\
    \addlinespace
    情報管理
      & \texttt{GetCurrentProcessId()}, \texttt{SetTimer()},
        \texttt{Sleep()}
      & \texttt{getpid()}, \texttt{alarm()}, \texttt{sleep()} \\
    \addlinespace
    通信
      & \texttt{CreatePipe()}, \texttt{CreateFileMapping()},
        \texttt{MapViewOfFile()}
      & \texttt{pipe()}, \texttt{shm\_open()}, \texttt{mmap()} \\
    \addlinespace
    保護
      & \texttt{SetFileSecurity()}, \texttt{InitializeSecurityDescriptor()},
        \texttt{SetSecurityDescriptorGroup()}
      & \texttt{chmod()}, \texttt{umask()}, \texttt{chown()} \\
    \bottomrule
  \end{tabular}
\end{table}

\section{オペレーティングシステムの構成}
\label{sec:l01-organization}

オペレーティングシステムは，そのサービスをどう構成するかによって異なる．
\source{P1L2，モノリシックOS，モジュール型OS，マイクロカーネル；
Silberschatz ら ch.~2．}\cref{fig:l01-structures} は3つの基本的な構成を
示し，\cref{tab:l01-organization} はそれらの利点と欠点をまとめたもので
ある．

\begin{figure}[htbp]
  \centering
  % Monolithic, modular and microkernel organization of an operating system.
% Usage: % Monolithic, modular and microkernel organization of an operating system.
% Usage: % Monolithic, modular and microkernel organization of an operating system.
% Usage: \input{figures/l01-structures}   (width ~116 mm)
\begin{tikzpicture}[x=1mm, y=1mm, font=\scriptsize\sffamily,
  app/.style={draw, rounded corners=1pt, fill=white, minimum height=5mm,
              minimum width=10mm, inner sep=0.5pt},
  row/.style={draw, fill=white, minimum height=4.6mm, minimum width=32mm,
              inner sep=0.5pt},
  mod/.style={draw, fill=white, minimum height=8mm, minimum width=10mm,
              inner sep=0.5pt, align=center},
  srv/.style={draw, rounded corners=1pt, fill=white, minimum height=8mm,
              minimum width=10.6mm, inner sep=0.5pt, align=center},
  kern/.style={draw=ln-accent, thick, fill=ln-box, rounded corners=2pt},
  hw/.style={draw, fill=white, minimum height=5mm, minimum width=36mm,
             inner sep=0.5pt},
  bnd/.style={dashed, ln-rule, thick},
  ttl/.style={font=\footnotesize\sffamily\bfseries, anchor=south},
  lab/.style={font=\scriptsize\sffamily, text=ln-muted},
  >={Stealth[length=1.4mm]}]
  % ================= monolithic
  \node[ttl] at (18,52) {\iflangja{モノリシック}{monolithic}};
  \foreach \x in {6,18,30} \node[app] at (\x,45) {\iflangja{アプリ}{app}};
  \draw[bnd] (0,39.5) -- (36,39.5);
  \draw[kern] (0,7) rectangle (36,38);
  \foreach \y/\t in {35.25/{\iflangja{プロセス・スレッド}{processes, threads}},
                     30.15/{\iflangja{スケジューリング}{scheduling}},
                     25.05/{\iflangja{メモリ管理}{memory management}},
                     19.95/{\iflangja{デバイスドライバ}{device drivers}},
                     14.85/{\iflangja{FS（逐次アクセス用）}{file system (sequential)}},
                     9.75/{\iflangja{FS（ランダムアクセス用）}{file system (random)}}} {
    \node[row] at (18,\y) {\t};
  }
  \node[hw] at (18,2.5) {\iflangja{ハードウェア}{hardware}};
  % ================= modular
  \begin{scope}[xshift=40mm]
  \node[ttl] at (18,52) {\iflangja{モジュール型}{modular}};
  \foreach \x in {6,18,30} \node[app] at (\x,45) {\iflangja{アプリ}{app}};
  \draw[bnd] (0,39.5) -- (36,39.5);
  \draw[kern] (0,7) rectangle (36,38);
  \node[row, minimum height=9mm, align=center] at (18,31.5)
    {\iflangja{基本サービス\\と API}{basic services\\and APIs}};
  \draw[very thick] (1.5,25.5) -- (34.5,25.5);
  \foreach \x/\t in {7/{FS},
                     18/{\iflangja{ドライバ}{driver}},
                     29/{\iflangja{プロト\\コル}{network\\protocol}}} {
    \node[mod] at (\x,19) {\t};
    \draw[very thick] (\x,23) -- (\x,25.5);
  }
  \node[lab, align=center] at (18,10.5)
    {\iflangja{モジュールを動的に追加}{modules loaded\\dynamically}};
  \node[hw] at (18,2.5) {\iflangja{ハードウェア}{hardware}};
  \end{scope}
  % ================= microkernel
  \begin{scope}[xshift=80mm]
  \node[ttl] at (18,52) {\iflangja{マイクロカーネル}{microkernel}};
  \foreach \x in {6,18,30} \node[app] at (\x,45) {\iflangja{アプリ}{app}};
  \node[srv] (fs) at (6,33)  {\iflangja{ファイル\\システム}{file\\system}};
  \node[srv] (dd) at (18,33) {\iflangja{デバイス\\ドライバ}{device\\driver}};
  \node[srv] (db) at (30,33) {\iflangja{データ\\ベース}{database}};
  \draw[bnd] (0,19.5) -- (36,19.5);
  \draw[kern] (0,7) rectangle (36,18);
  \node[align=center] at (18,12.5)
    {\iflangja{アドレス空間，\\スレッド，IPC}{address spaces,\\threads, IPC}};
  \foreach \n in {fs,dd,db} \draw[<->, ln-accent] (\n.south) -- ++(0,-10);
  \node[hw] at (18,2.5) {\iflangja{ハードウェア}{hardware}};
  \end{scope}
\end{tikzpicture}
   (width ~116 mm)
\begin{tikzpicture}[x=1mm, y=1mm, font=\scriptsize\sffamily,
  app/.style={draw, rounded corners=1pt, fill=white, minimum height=5mm,
              minimum width=10mm, inner sep=0.5pt},
  row/.style={draw, fill=white, minimum height=4.6mm, minimum width=32mm,
              inner sep=0.5pt},
  mod/.style={draw, fill=white, minimum height=8mm, minimum width=10mm,
              inner sep=0.5pt, align=center},
  srv/.style={draw, rounded corners=1pt, fill=white, minimum height=8mm,
              minimum width=10.6mm, inner sep=0.5pt, align=center},
  kern/.style={draw=ln-accent, thick, fill=ln-box, rounded corners=2pt},
  hw/.style={draw, fill=white, minimum height=5mm, minimum width=36mm,
             inner sep=0.5pt},
  bnd/.style={dashed, ln-rule, thick},
  ttl/.style={font=\footnotesize\sffamily\bfseries, anchor=south},
  lab/.style={font=\scriptsize\sffamily, text=ln-muted},
  >={Stealth[length=1.4mm]}]
  % ================= monolithic
  \node[ttl] at (18,52) {\iflangja{モノリシック}{monolithic}};
  \foreach \x in {6,18,30} \node[app] at (\x,45) {\iflangja{アプリ}{app}};
  \draw[bnd] (0,39.5) -- (36,39.5);
  \draw[kern] (0,7) rectangle (36,38);
  \foreach \y/\t in {35.25/{\iflangja{プロセス・スレッド}{processes, threads}},
                     30.15/{\iflangja{スケジューリング}{scheduling}},
                     25.05/{\iflangja{メモリ管理}{memory management}},
                     19.95/{\iflangja{デバイスドライバ}{device drivers}},
                     14.85/{\iflangja{FS（逐次アクセス用）}{file system (sequential)}},
                     9.75/{\iflangja{FS（ランダムアクセス用）}{file system (random)}}} {
    \node[row] at (18,\y) {\t};
  }
  \node[hw] at (18,2.5) {\iflangja{ハードウェア}{hardware}};
  % ================= modular
  \begin{scope}[xshift=40mm]
  \node[ttl] at (18,52) {\iflangja{モジュール型}{modular}};
  \foreach \x in {6,18,30} \node[app] at (\x,45) {\iflangja{アプリ}{app}};
  \draw[bnd] (0,39.5) -- (36,39.5);
  \draw[kern] (0,7) rectangle (36,38);
  \node[row, minimum height=9mm, align=center] at (18,31.5)
    {\iflangja{基本サービス\\と API}{basic services\\and APIs}};
  \draw[very thick] (1.5,25.5) -- (34.5,25.5);
  \foreach \x/\t in {7/{FS},
                     18/{\iflangja{ドライバ}{driver}},
                     29/{\iflangja{プロト\\コル}{network\\protocol}}} {
    \node[mod] at (\x,19) {\t};
    \draw[very thick] (\x,23) -- (\x,25.5);
  }
  \node[lab, align=center] at (18,10.5)
    {\iflangja{モジュールを動的に追加}{modules loaded\\dynamically}};
  \node[hw] at (18,2.5) {\iflangja{ハードウェア}{hardware}};
  \end{scope}
  % ================= microkernel
  \begin{scope}[xshift=80mm]
  \node[ttl] at (18,52) {\iflangja{マイクロカーネル}{microkernel}};
  \foreach \x in {6,18,30} \node[app] at (\x,45) {\iflangja{アプリ}{app}};
  \node[srv] (fs) at (6,33)  {\iflangja{ファイル\\システム}{file\\system}};
  \node[srv] (dd) at (18,33) {\iflangja{デバイス\\ドライバ}{device\\driver}};
  \node[srv] (db) at (30,33) {\iflangja{データ\\ベース}{database}};
  \draw[bnd] (0,19.5) -- (36,19.5);
  \draw[kern] (0,7) rectangle (36,18);
  \node[align=center] at (18,12.5)
    {\iflangja{アドレス空間，\\スレッド，IPC}{address spaces,\\threads, IPC}};
  \foreach \n in {fs,dd,db} \draw[<->, ln-accent] (\n.south) -- ++(0,-10);
  \node[hw] at (18,2.5) {\iflangja{ハードウェア}{hardware}};
  \end{scope}
\end{tikzpicture}
   (width ~116 mm)
\begin{tikzpicture}[x=1mm, y=1mm, font=\scriptsize\sffamily,
  app/.style={draw, rounded corners=1pt, fill=white, minimum height=5mm,
              minimum width=10mm, inner sep=0.5pt},
  row/.style={draw, fill=white, minimum height=4.6mm, minimum width=32mm,
              inner sep=0.5pt},
  mod/.style={draw, fill=white, minimum height=8mm, minimum width=10mm,
              inner sep=0.5pt, align=center},
  srv/.style={draw, rounded corners=1pt, fill=white, minimum height=8mm,
              minimum width=10.6mm, inner sep=0.5pt, align=center},
  kern/.style={draw=ln-accent, thick, fill=ln-box, rounded corners=2pt},
  hw/.style={draw, fill=white, minimum height=5mm, minimum width=36mm,
             inner sep=0.5pt},
  bnd/.style={dashed, ln-rule, thick},
  ttl/.style={font=\footnotesize\sffamily\bfseries, anchor=south},
  lab/.style={font=\scriptsize\sffamily, text=ln-muted},
  >={Stealth[length=1.4mm]}]
  % ================= monolithic
  \node[ttl] at (18,52) {\iflangja{モノリシック}{monolithic}};
  \foreach \x in {6,18,30} \node[app] at (\x,45) {\iflangja{アプリ}{app}};
  \draw[bnd] (0,39.5) -- (36,39.5);
  \draw[kern] (0,7) rectangle (36,38);
  \foreach \y/\t in {35.25/{\iflangja{プロセス・スレッド}{processes, threads}},
                     30.15/{\iflangja{スケジューリング}{scheduling}},
                     25.05/{\iflangja{メモリ管理}{memory management}},
                     19.95/{\iflangja{デバイスドライバ}{device drivers}},
                     14.85/{\iflangja{FS（逐次アクセス用）}{file system (sequential)}},
                     9.75/{\iflangja{FS（ランダムアクセス用）}{file system (random)}}} {
    \node[row] at (18,\y) {\t};
  }
  \node[hw] at (18,2.5) {\iflangja{ハードウェア}{hardware}};
  % ================= modular
  \begin{scope}[xshift=40mm]
  \node[ttl] at (18,52) {\iflangja{モジュール型}{modular}};
  \foreach \x in {6,18,30} \node[app] at (\x,45) {\iflangja{アプリ}{app}};
  \draw[bnd] (0,39.5) -- (36,39.5);
  \draw[kern] (0,7) rectangle (36,38);
  \node[row, minimum height=9mm, align=center] at (18,31.5)
    {\iflangja{基本サービス\\と API}{basic services\\and APIs}};
  \draw[very thick] (1.5,25.5) -- (34.5,25.5);
  \foreach \x/\t in {7/{FS},
                     18/{\iflangja{ドライバ}{driver}},
                     29/{\iflangja{プロト\\コル}{network\\protocol}}} {
    \node[mod] at (\x,19) {\t};
    \draw[very thick] (\x,23) -- (\x,25.5);
  }
  \node[lab, align=center] at (18,10.5)
    {\iflangja{モジュールを動的に追加}{modules loaded\\dynamically}};
  \node[hw] at (18,2.5) {\iflangja{ハードウェア}{hardware}};
  \end{scope}
  % ================= microkernel
  \begin{scope}[xshift=80mm]
  \node[ttl] at (18,52) {\iflangja{マイクロカーネル}{microkernel}};
  \foreach \x in {6,18,30} \node[app] at (\x,45) {\iflangja{アプリ}{app}};
  \node[srv] (fs) at (6,33)  {\iflangja{ファイル\\システム}{file\\system}};
  \node[srv] (dd) at (18,33) {\iflangja{デバイス\\ドライバ}{device\\driver}};
  \node[srv] (db) at (30,33) {\iflangja{データ\\ベース}{database}};
  \draw[bnd] (0,19.5) -- (36,19.5);
  \draw[kern] (0,7) rectangle (36,18);
  \node[align=center] at (18,12.5)
    {\iflangja{アドレス空間，\\スレッド，IPC}{address spaces,\\threads, IPC}};
  \foreach \n in {fs,dd,db} \draw[<->, ln-accent] (\n.south) -- ++(0,-10);
  \node[hw] at (18,2.5) {\iflangja{ハードウェア}{hardware}};
  \end{scope}
\end{tikzpicture}

  \caption{オペレーティングシステムの3つの構成．破線はユーザ/カーネル
    境界である．モノリシックカーネルはすべてのサービスを含み，
    モジュール型カーネルは共通のインタフェースを通じてサービスを
    モジュールとして追加し，マイクロカーネルはアドレス空間，スレッド，
    IPC だけを保持して他のすべてのサービスをユーザレベルで実行する．}
  \label{fig:l01-structures}
\end{figure}

\subsection{モノリシックOS}

歴史的には，オペレーティングシステムはモノリシックであった．
\term[ものりしっくOS]{モノリシックOS}[monolithic operating system]では，
どのアプリケーションが必要としうるサービスも，どの種類のハードウェアが
要求しうるサービスも，すべてオペレーティングシステムの一部である．
メモリ管理，デバイスドライバ，ファイル管理，プロセスとスレッド，
スケジューリングがそうである．このようなシステムは，逐次アクセスに特化
したものと，データベースが使うようなランダムアクセスに最適化したものの
ように，複数のファイルシステムを含むことさえある．これらのコードは
すべて一緒にコンパイルされるので，コンパイラは構成要素の境界を越えて，
例えばインライン化によって最適化できる．その代償は，カスタマイズ，移植，
保守，デバッグ，更新が難しい大量のコードと状態，そしてアプリケーションに
残されるメモリを減らし，その性能を低下させうる大きなメモリフットプリント
である．

\subsection{モジュール型OS}

Linux を標準的な例とする\term[もじゅーるがたOS]{モジュール型OS}[modular operating system]%
は，基本的なサービスと API の集合を備え，それ以外のものは
\term[かーねるもじゅーる]{カーネルモジュール}[kernel module]として追加
できる．オペレーティングシステムは，モジュールがその一部となるために
実装しなければならないインタフェースを，例えばファイルシステム用の
インタフェースとして定める．モジュールはワークロードに応じて動的に
インストールできる．データベースを実行するマシンは，ランダムアクセスに
最適化されたファイルシステムをロードできる．必要なモジュールだけがロード
されるので，システムはより小さく，保守と更新が容易で，アプリケーションに
より多くの資源を残す．モジュールインタフェースを介する間接参照は多少の
性能低下を招きうるが，通常はわずかである．一方，モジュールはまったく異なる
コードベースに由来しうるので，バグの発生源となる．\margin{Silberschatz
らは Linux を，ロード可能なモジュールを持つモノリシックカーネルとして
説明している．}

\subsection{マイクロカーネル}

\term[まいくろかーねる]{マイクロカーネル}[microkernel]は，オペレーティング
システムのレベルで最も基本的なプリミティブ，すなわち実行中の
アプリケーションのメモリを表すアドレス空間と，その実行コンテキストを表す
スレッドだけを提供する．それ以外のすべて---ファイルシステム，デバイス
ドライバ，データベース---は，カーネルの外，ユーザレベルで別々のプロセス
として実行される．したがって1つの要求を処理するにはプロセス間の多数の
やり取りが必要であり，プロセス間通信（IPC）は，アドレス空間とスレッドと
並ぶマイクロカーネルの中核的な抽象化かつ機構である．マイクロカーネルは
小さいので，カーネル自体のメモリフットプリントは小さく，アプリケーション
によってはオーバーヘッドが低く性能もよい．そのコードは検証もしやすい．
そのため，オペレーティングシステムの正しい振る舞いが決定的に重要な
組込みデバイスや制御システムでは，マイクロカーネルが魅力的である．他方で
マイクロカーネルは下位のハードウェアに特化していることが多く，移植性が
制限される．また共通のオペレーティングシステム構成要素を見つけにくいので
ソフトウェア開発はより複雑になり，ユーザ/カーネル境界を頻繁に越えることは
コストが高い．

\begin{table}[htbp]
  \centering
  \caption{3つの構成の利点と欠点．}
  \label{tab:l01-organization}
  \small
  \begin{tabular}{@{}>{\RaggedRight}p{19mm}>{\RaggedRight}p{46mm}%
      >{\RaggedRight\arraybackslash}p{47mm}@{}}
    \toprule
    構成 & 利点 & 欠点 \\
    \midrule
    モノリシック & すべてを含む．構成要素を越えた最適化ができる
      & カスタマイズ，移植，管理が難しい．メモリフットプリントが大きい \\
    \addlinespace
    モジュール型 & 保守しやすい．フットプリントが小さい．必要な資源が
      少ない
      & インタフェースを介する間接参照．別々に開発されたモジュールに
      起因するバグ \\
    \addlinespace
    マイクロカーネル & カーネルのフットプリントが小さい．検証しやすい
      & 移植性が限られる．ソフトウェア開発が複雑．ユーザ/カーネル境界を
      頻繁に越える \\
    \bottomrule
  \end{tabular}
\end{table}

\begin{example}
  ディスク上のファイルからの読み出しを考える．モノリシックカーネルや
  モジュール型カーネルではファイルシステムとディスクドライバがカーネル
  内で実行されるので，アプリケーションは1回のシステムコールを行い，
  要求は入るときと戻るときの2回，ユーザ/カーネル境界を越える．
  マイクロカーネルでは，アプリケーションがファイルシステムプロセスへ
  IPC メッセージを送り，そのプロセスがドライバプロセスへメッセージを
  送る．ドライバはファイルシステムへ返答し，ファイルシステムは
  アプリケーションへ返答する．4つのメッセージはそれぞれカーネルに入り，
  また出ていくので，合計8回の横断となる．1回の横断を
  \SI{80}{\nano\second}，読み出しを毎秒 \num{50000} 回とすると，
  モノリシックな設計は1秒当たり
  $\num{50000}\times 2\times\SI{80}{\nano\second} =
  \SI{8}{\milli\second}$ を，マイクロカーネルは
  $\num{50000}\times 8\times\SI{80}{\nano\second} =
  \SI{32}{\milli\second}$ を横断に費やす．これには，関係するプロセスの
  アドレス空間の切り替えは含まれていない．
\end{example}

\begin{keypoint}
  モノリシックカーネルは完全で全体として最適化できるが，大きく管理が
  難しい．モジュール型カーネルは必要なサービスだけを追加するが，間接参照
  の代償を払う．マイクロカーネルは小さく検証しやすいが，頻繁なユーザ/
  カーネル境界の横断と IPC の代償を払う．
\end{keypoint}

\section{Linux と Mac OS X}
\label{sec:l01-linux-macos}

Linux システムは階層状に構成される（\cref{fig:l01-linux} 左）．
\source{P1L2，Linux と Mac OS のアーキテクチャ；Tanenbaum と Bos
\cite{tanenbaum2015} ch.~10；Silberschatz ら ch.~2．}Linux カーネルは
ハードウェア上で直接カーネルモードで実行される．その上位のユーザモードでは，
標準ライブラリが \texttt{open}，\texttt{close}，\texttt{read}，
\texttt{write}，\texttt{fork} などの関数としてシステムコールインタフェース
を提供し，シェル，エディタ，コンパイラなどの標準ユーティリティプログラムが
そのライブラリを使い，利用者はユーティリティプログラムや自分の
アプリケーションとやり取りする．これらの層の間の境界が，システム
コールインタフェース，ライブラリインタフェース，ユーザインタフェース
である．

\begin{figure}[htbp]
  \centering
  % Linux: the layers of a Linux system (left) and the structure of the kernel (right).
% Usage: % Linux: the layers of a Linux system (left) and the structure of the kernel (right).
% Usage: % Linux: the layers of a Linux system (left) and the structure of the kernel (right).
% Usage: \input{figures/l01-linux}   (width ~118 mm)
\begin{tikzpicture}[x=1mm, y=1mm, font=\scriptsize\sffamily,
  lay/.style={draw, rounded corners=1pt, fill=white, minimum width=40mm,
              text width=38mm, align=center, inner sep=0.5pt},
  bar/.style={draw, fill=ln-box, minimum width=72mm, minimum height=5mm,
              inner sep=0.5pt},
  hd/.style={draw, fill=ln-box, minimum width=23.2mm, minimum height=5mm,
             inner sep=0.5pt, font=\scriptsize\sffamily\bfseries},
  cell/.style={draw, fill=white, minimum width=23.2mm, text width=22.2mm,
               minimum height=9mm, align=center, inner sep=0.5pt},
  lab/.style={font=\scriptsize\sffamily, text=ln-muted},
  ttl/.style={font=\footnotesize\sffamily\bfseries, anchor=base}]
  % ---------------- layers
  \node[ttl] at (20,54) {\iflangja{Linux システムの階層}{layers of a Linux system}};
  \node[lay, minimum height=6mm] at (20,48.5) {\iflangja{ユーザ}{users}};
  \node[lay, minimum height=10mm] at (20,39)
    {\iflangja{標準ユーティリティ\\（シェル，エディタ，コンパイラ，…）}{standard utility programs\\(shell, editors, compilers, \dots)}};
  \node[lay, minimum height=10mm] at (20,27.5)
    {\iflangja{標準ライブラリ\\（open, close, read, write, fork, …）}{standard library\\(open, close, read, write, fork, \dots)}};
  \node[lay, minimum height=8mm, fill=ln-box, draw=ln-accent, thick] at (20,15.5)
    {\iflangja{Linux カーネル}{Linux kernel}};
  \node[lay, minimum height=8mm] at (20,5.5)
    {\iflangja{ハードウェア}{hardware}};
  \draw[dashed, ln-rule, thick] (-1,21) -- (41,21);
  % ---------------- kernel structure
  \begin{scope}[xshift=46mm]
  \node[ttl] at (36,54) {\iflangja{Linux カーネルの構成}{structure of the Linux kernel}};
  \node[bar] at (36,46) {\iflangja{システムコール}{system calls}};
  \node[hd] at (11.6,39.5) {I/O};
  \node[hd] at (36,39.5)   {\iflangja{メモリ管理}{memory mgmt.}};
  \node[hd] at (60.4,39.5) {\iflangja{プロセス管理}{process mgmt.}};
  \foreach \y/\a/\b/\c in {
      31.5/{\iflangja{端末\\文字デバイスドライバ}{terminals, character\\device drivers}}/{\iflangja{仮想メモリ}{virtual memory}}/{\iflangja{シグナル処理}{signal handling}},
      21.5/{\iflangja{ソケット，プロトコル，\\ネットワークドライバ}{sockets, protocols,\\network drivers}}/{\iflangja{ページング，\\ページ置換}{paging, page\\replacement}}/{\iflangja{プロセス・スレッドの\\生成と終了}{process and thread\\creation, termination}},
      11.5/{\iflangja{ファイルシステム，\\ブロックデバイスドライバ}{file systems, block\\device drivers}}/{\iflangja{ページキャッシュ}{page cache}}/{\iflangja{CPU スケジューリング}{CPU scheduling}}} {
    \node[cell] at (11.6,\y) {\a};
    \node[cell] at (36,\y)   {\b};
    \node[cell] at (60.4,\y) {\c};
  }
  \node[bar] at (36,3.5) {\iflangja{割り込み，ディスパッチャ}{interrupts, dispatcher}};
  \end{scope}
\end{tikzpicture}
   (width ~118 mm)
\begin{tikzpicture}[x=1mm, y=1mm, font=\scriptsize\sffamily,
  lay/.style={draw, rounded corners=1pt, fill=white, minimum width=40mm,
              text width=38mm, align=center, inner sep=0.5pt},
  bar/.style={draw, fill=ln-box, minimum width=72mm, minimum height=5mm,
              inner sep=0.5pt},
  hd/.style={draw, fill=ln-box, minimum width=23.2mm, minimum height=5mm,
             inner sep=0.5pt, font=\scriptsize\sffamily\bfseries},
  cell/.style={draw, fill=white, minimum width=23.2mm, text width=22.2mm,
               minimum height=9mm, align=center, inner sep=0.5pt},
  lab/.style={font=\scriptsize\sffamily, text=ln-muted},
  ttl/.style={font=\footnotesize\sffamily\bfseries, anchor=base}]
  % ---------------- layers
  \node[ttl] at (20,54) {\iflangja{Linux システムの階層}{layers of a Linux system}};
  \node[lay, minimum height=6mm] at (20,48.5) {\iflangja{ユーザ}{users}};
  \node[lay, minimum height=10mm] at (20,39)
    {\iflangja{標準ユーティリティ\\（シェル，エディタ，コンパイラ，…）}{standard utility programs\\(shell, editors, compilers, \dots)}};
  \node[lay, minimum height=10mm] at (20,27.5)
    {\iflangja{標準ライブラリ\\（open, close, read, write, fork, …）}{standard library\\(open, close, read, write, fork, \dots)}};
  \node[lay, minimum height=8mm, fill=ln-box, draw=ln-accent, thick] at (20,15.5)
    {\iflangja{Linux カーネル}{Linux kernel}};
  \node[lay, minimum height=8mm] at (20,5.5)
    {\iflangja{ハードウェア}{hardware}};
  \draw[dashed, ln-rule, thick] (-1,21) -- (41,21);
  % ---------------- kernel structure
  \begin{scope}[xshift=46mm]
  \node[ttl] at (36,54) {\iflangja{Linux カーネルの構成}{structure of the Linux kernel}};
  \node[bar] at (36,46) {\iflangja{システムコール}{system calls}};
  \node[hd] at (11.6,39.5) {I/O};
  \node[hd] at (36,39.5)   {\iflangja{メモリ管理}{memory mgmt.}};
  \node[hd] at (60.4,39.5) {\iflangja{プロセス管理}{process mgmt.}};
  \foreach \y/\a/\b/\c in {
      31.5/{\iflangja{端末\\文字デバイスドライバ}{terminals, character\\device drivers}}/{\iflangja{仮想メモリ}{virtual memory}}/{\iflangja{シグナル処理}{signal handling}},
      21.5/{\iflangja{ソケット，プロトコル，\\ネットワークドライバ}{sockets, protocols,\\network drivers}}/{\iflangja{ページング，\\ページ置換}{paging, page\\replacement}}/{\iflangja{プロセス・スレッドの\\生成と終了}{process and thread\\creation, termination}},
      11.5/{\iflangja{ファイルシステム，\\ブロックデバイスドライバ}{file systems, block\\device drivers}}/{\iflangja{ページキャッシュ}{page cache}}/{\iflangja{CPU スケジューリング}{CPU scheduling}}} {
    \node[cell] at (11.6,\y) {\a};
    \node[cell] at (36,\y)   {\b};
    \node[cell] at (60.4,\y) {\c};
  }
  \node[bar] at (36,3.5) {\iflangja{割り込み，ディスパッチャ}{interrupts, dispatcher}};
  \end{scope}
\end{tikzpicture}
   (width ~118 mm)
\begin{tikzpicture}[x=1mm, y=1mm, font=\scriptsize\sffamily,
  lay/.style={draw, rounded corners=1pt, fill=white, minimum width=40mm,
              text width=38mm, align=center, inner sep=0.5pt},
  bar/.style={draw, fill=ln-box, minimum width=72mm, minimum height=5mm,
              inner sep=0.5pt},
  hd/.style={draw, fill=ln-box, minimum width=23.2mm, minimum height=5mm,
             inner sep=0.5pt, font=\scriptsize\sffamily\bfseries},
  cell/.style={draw, fill=white, minimum width=23.2mm, text width=22.2mm,
               minimum height=9mm, align=center, inner sep=0.5pt},
  lab/.style={font=\scriptsize\sffamily, text=ln-muted},
  ttl/.style={font=\footnotesize\sffamily\bfseries, anchor=base}]
  % ---------------- layers
  \node[ttl] at (20,54) {\iflangja{Linux システムの階層}{layers of a Linux system}};
  \node[lay, minimum height=6mm] at (20,48.5) {\iflangja{ユーザ}{users}};
  \node[lay, minimum height=10mm] at (20,39)
    {\iflangja{標準ユーティリティ\\（シェル，エディタ，コンパイラ，…）}{standard utility programs\\(shell, editors, compilers, \dots)}};
  \node[lay, minimum height=10mm] at (20,27.5)
    {\iflangja{標準ライブラリ\\（open, close, read, write, fork, …）}{standard library\\(open, close, read, write, fork, \dots)}};
  \node[lay, minimum height=8mm, fill=ln-box, draw=ln-accent, thick] at (20,15.5)
    {\iflangja{Linux カーネル}{Linux kernel}};
  \node[lay, minimum height=8mm] at (20,5.5)
    {\iflangja{ハードウェア}{hardware}};
  \draw[dashed, ln-rule, thick] (-1,21) -- (41,21);
  % ---------------- kernel structure
  \begin{scope}[xshift=46mm]
  \node[ttl] at (36,54) {\iflangja{Linux カーネルの構成}{structure of the Linux kernel}};
  \node[bar] at (36,46) {\iflangja{システムコール}{system calls}};
  \node[hd] at (11.6,39.5) {I/O};
  \node[hd] at (36,39.5)   {\iflangja{メモリ管理}{memory mgmt.}};
  \node[hd] at (60.4,39.5) {\iflangja{プロセス管理}{process mgmt.}};
  \foreach \y/\a/\b/\c in {
      31.5/{\iflangja{端末\\文字デバイスドライバ}{terminals, character\\device drivers}}/{\iflangja{仮想メモリ}{virtual memory}}/{\iflangja{シグナル処理}{signal handling}},
      21.5/{\iflangja{ソケット，プロトコル，\\ネットワークドライバ}{sockets, protocols,\\network drivers}}/{\iflangja{ページング，\\ページ置換}{paging, page\\replacement}}/{\iflangja{プロセス・スレッドの\\生成と終了}{process and thread\\creation, termination}},
      11.5/{\iflangja{ファイルシステム，\\ブロックデバイスドライバ}{file systems, block\\device drivers}}/{\iflangja{ページキャッシュ}{page cache}}/{\iflangja{CPU スケジューリング}{CPU scheduling}}} {
    \node[cell] at (11.6,\y) {\a};
    \node[cell] at (36,\y)   {\b};
    \node[cell] at (60.4,\y) {\c};
  }
  \node[bar] at (36,3.5) {\iflangja{割り込み，ディスパッチャ}{interrupts, dispatcher}};
  \end{scope}
\end{tikzpicture}

  \caption{左: Linux システムの階層．破線はユーザ/カーネル境界である．
    右: Linux カーネルの構成要素（Tanenbaum と Bos \cite{tanenbaum2015}
    による）．}
  \label{fig:l01-linux}
\end{figure}

カーネル自体（\cref{fig:l01-linux} 右）へは上部のシステムコールから
入る．下部では割り込みを処理し，CPU をプロセス間で切り替える
ディスパッチャを備える．その間に3つの構成要素がある．I/O の構成要素は，
端末と文字デバイスのドライバ，ソケット，ネットワークプロトコルと
ネットワークデバイスのドライバ，ファイルシステムとブロックデバイスの
ドライバを含む．メモリ管理の構成要素は，仮想メモリ，ページングと
ページ置換，ページキャッシュを実装する．プロセス管理の構成要素は，
シグナルを扱い，プロセスとスレッドを生成・終了し，CPU をスケジュール
する．各構成要素は他から独立に変更・置換でき，これがモジュール型の
方式を可能にしている．

Mac OS X（現在の macOS）は Mach マイクロカーネルを中心に構築されている
（\cref{fig:l01-macos}）．\margin{Mach，BSD，I/O Kit，カーネル拡張は
1つのカーネルアドレス空間を共有する．すなわちハイブリッドな設計である
（Silberschatz ら）．}Mach は，メモリ管理，スレッドスケジューリング，
そして遠隔手続き呼び出し（第13章）を含む IPC という主要なプリミティブを
実装する．BSD の構成要素は，BSD のコマンドラインインタフェース，POSIX
API のサポート，ネットワーク I/O によって Unix との相互運用性を提供する．
すべてのアプリケーション環境はこれら2つの上に層をなす．その下では，
I/O Kit がデバイスドライバを開発するための環境であり，カーネル拡張は
動的にカーネルへロードされる構成要素である．

\begin{figure}[htbp]
  \centering
  % Structure of Mac OS X: Mach microkernel with a BSD layer, I/O Kit and kernel extensions.
% I/O Kit and kernel extensions run in kernel mode, so they lie inside the
% kernel environment.
% Usage: % Structure of Mac OS X: Mach microkernel with a BSD layer, I/O Kit and kernel extensions.
% I/O Kit and kernel extensions run in kernel mode, so they lie inside the
% kernel environment.
% Usage: % Structure of Mac OS X: Mach microkernel with a BSD layer, I/O Kit and kernel extensions.
% I/O Kit and kernel extensions run in kernel mode, so they lie inside the
% kernel environment.
% Usage: \input{figures/l01-macos}   (width ~96 mm)
\begin{tikzpicture}[x=1mm, y=1mm, font=\scriptsize\sffamily,
  blk/.style={draw, rounded corners=1pt, fill=white, align=center, inner sep=1pt},
  lab/.style={font=\scriptsize\sffamily\bfseries, text=ln-accent}]
  \node[blk, minimum width=96mm, minimum height=7mm] at (48,43)
    {\iflangja{アプリケーション環境とサービス}{application environments and services}};
  \draw[draw=ln-accent, thick, fill=ln-box, rounded corners=2pt] (0,1) rectangle (96,38);
  \node[lab, anchor=north west] at (1,37.5) {\iflangja{カーネル環境}{kernel environment}};
  \node[blk, minimum width=92mm, minimum height=8mm] at (48,26)
    {\textbf{BSD}\quad \iflangja{BSD コマンドライン，POSIX API，ネットワーク I/O}{BSD command line, POSIX API, network I/O}};
  \node[blk, minimum width=92mm, minimum height=8mm] at (48,16)
    {\textbf{Mach}\quad \iflangja{メモリ管理，スレッドスケジューリング，IPC・RPC}{memory management, thread scheduling, IPC and RPC}};
  \node[blk, minimum width=45.5mm, minimum height=7mm] at (24.75,6.5)
    {\textbf{I/O Kit}\quad \iflangja{デバイスドライバ}{device drivers}};
  \node[blk, minimum width=45.5mm, minimum height=7mm] at (71.25,6.5)
    {\iflangja{カーネル拡張}{kernel extensions}};
\end{tikzpicture}
   (width ~96 mm)
\begin{tikzpicture}[x=1mm, y=1mm, font=\scriptsize\sffamily,
  blk/.style={draw, rounded corners=1pt, fill=white, align=center, inner sep=1pt},
  lab/.style={font=\scriptsize\sffamily\bfseries, text=ln-accent}]
  \node[blk, minimum width=96mm, minimum height=7mm] at (48,43)
    {\iflangja{アプリケーション環境とサービス}{application environments and services}};
  \draw[draw=ln-accent, thick, fill=ln-box, rounded corners=2pt] (0,1) rectangle (96,38);
  \node[lab, anchor=north west] at (1,37.5) {\iflangja{カーネル環境}{kernel environment}};
  \node[blk, minimum width=92mm, minimum height=8mm] at (48,26)
    {\textbf{BSD}\quad \iflangja{BSD コマンドライン，POSIX API，ネットワーク I/O}{BSD command line, POSIX API, network I/O}};
  \node[blk, minimum width=92mm, minimum height=8mm] at (48,16)
    {\textbf{Mach}\quad \iflangja{メモリ管理，スレッドスケジューリング，IPC・RPC}{memory management, thread scheduling, IPC and RPC}};
  \node[blk, minimum width=45.5mm, minimum height=7mm] at (24.75,6.5)
    {\textbf{I/O Kit}\quad \iflangja{デバイスドライバ}{device drivers}};
  \node[blk, minimum width=45.5mm, minimum height=7mm] at (71.25,6.5)
    {\iflangja{カーネル拡張}{kernel extensions}};
\end{tikzpicture}
   (width ~96 mm)
\begin{tikzpicture}[x=1mm, y=1mm, font=\scriptsize\sffamily,
  blk/.style={draw, rounded corners=1pt, fill=white, align=center, inner sep=1pt},
  lab/.style={font=\scriptsize\sffamily\bfseries, text=ln-accent}]
  \node[blk, minimum width=96mm, minimum height=7mm] at (48,43)
    {\iflangja{アプリケーション環境とサービス}{application environments and services}};
  \draw[draw=ln-accent, thick, fill=ln-box, rounded corners=2pt] (0,1) rectangle (96,38);
  \node[lab, anchor=north west] at (1,37.5) {\iflangja{カーネル環境}{kernel environment}};
  \node[blk, minimum width=92mm, minimum height=8mm] at (48,26)
    {\textbf{BSD}\quad \iflangja{BSD コマンドライン，POSIX API，ネットワーク I/O}{BSD command line, POSIX API, network I/O}};
  \node[blk, minimum width=92mm, minimum height=8mm] at (48,16)
    {\textbf{Mach}\quad \iflangja{メモリ管理，スレッドスケジューリング，IPC・RPC}{memory management, thread scheduling, IPC and RPC}};
  \node[blk, minimum width=45.5mm, minimum height=7mm] at (24.75,6.5)
    {\textbf{I/O Kit}\quad \iflangja{デバイスドライバ}{device drivers}};
  \node[blk, minimum width=45.5mm, minimum height=7mm] at (71.25,6.5)
    {\iflangja{カーネル拡張}{kernel extensions}};
\end{tikzpicture}

  \caption{Mac OS X の構造．BSD 層，I/O Kit，カーネル拡張を伴う Mach
    マイクロカーネルであり，これらはすべてカーネルモードで実行される．}
  \label{fig:l01-macos}
\end{figure}

\begin{keypoint}[章のまとめ]
  オペレーティングシステムは，アプリケーションとハードウェアの間にある
  特権的なソフトウェアの層である．ハードウェアの複雑さを抽象化の背後に
  隠し，ハードウェア資源の利用を調停し，アプリケーションを互いに隔離し
  保護する．それは抽象化，機構，方針から構築され，機構と方針の分離と，
  頻繁に起こる場合の最適化がその設計を導く．ハードウェアに支えられた
  ユーザモードとカーネルモードがオペレーティングシステムを保護し，
  アプリケーションは安価ではないシステムコールとトラップによって境界を
  越え，オペレーティングシステムはシグナルでアプリケーションに通知する．
  モノリシック，モジュール型，マイクロカーネルの構成は，完全性，柔軟性，
  規模，性能を互いにトレードオフしており，モジュール型の Linux カーネル
  と Mach に基づく Mac OS X がそれを例示している．
\end{keypoint}

\end{document}
